\documentclass[11pt]{article}
\usepackage[margin=1in]{geometry}
\usepackage{amsmath,amssymb}
\usepackage{booktabs}
\usepackage{siunitx}
\usepackage[hidelinks]{hyperref}
\usepackage{xcolor}
\usepackage{tcolorbox}

\title{Replication Report:\\
\emph{Ferroelectric control of magnetic skyrmions in two-dimensional\\
van der Waals heterostructures}\\
\large Huang, Shao \& Tsymbal --- arXiv:2202.11348}
\author{REPLICATE-PROJECT / TEXTURES-100 (textures-polar-huang2022)}
\date{Independent reimplementation --- no author code}

\begin{document}
\maketitle

\begin{abstract}
We independently reimplement the core physics of Huang \emph{et al.}
(arXiv:2202.11348): ferroelectric (FE) control of magnetic skyrmions in
Fe$_3$GeTe$_2$/In$_2$Se$_3$ van der Waals heterostructures. The paper predicts,
from DFT + atomistic spin-dynamics, that reversing the In$_2$Se$_3$ FE
polarization switches the interfacial Dzyaloshinskii--Moriya interaction (DMI)
between $D_\uparrow=0.28$ and $D_\downarrow=0.06$~mJ/m$^2$, thereby
\emph{creating} vs.\ \emph{annihilating} a $\sim$12~nm N\'eel skyrmion, and that
in the trilayer the DMI sign flips ($D_{\uparrow\uparrow}\approx0.22$,
$D_{\downarrow\downarrow}\approx-0.24$~mJ/m$^2$) reversing skyrmion chirality.
We rebuild a discrete square-lattice micromagnetic model (exchange + interfacial
N\'eel DMI + uniaxial anisotropy) and relax a circular seed by overdamped LLG.
The analytic critical DMI, $D_c=(4/\pi)\sqrt{AK}=0.255$~mJ/m$^2$, brackets the
paper's cases exactly: $D_\uparrow=0.28>D_c>D_\downarrow=0.06$, reproducing the
create/annihilate switch with zero fitting. Topological charge $|Q|\approx1$
confirms genuine skyrmions; larger $|D|$ yields larger skyrmions (trend
reproduced).

\begin{tcolorbox}[colback=gray!8,colframe=black]
\textbf{VERDICT: PARTIAL} (leaning REPLICATED on mechanism).\quad
Coverage $\approx 6$--$7/10$, \; Agreement $\approx 7/10$.\\
Core mechanism (DMI-controlled skyrmion creation/annihilation via FE switching,
and chirality tied to sign$(D)$) is \textbf{fully reproduced}. Absolute
skyrmion diameter is \textbf{off by $\sim$2--3$\times$} because the exchange
stiffness $A$ appears only in the paper's Fig.~4 (not in extractable text);
diameter $\propto\sqrt{A}$.
\end{tcolorbox}
\end{abstract}

\section{The paper's claims}
\begin{enumerate}
\item \textbf{FE-switchable DMI (bilayer).} DFT gives $D_\uparrow=0.28$~mJ/m$^2$
(polarization up) and $D_\downarrow=0.06$~mJ/m$^2$ (down).
\item \textbf{Create / annihilate switch.} With adopted anisotropy
$K\approx0.04$~MJ/m$^3$, a $\sim$12~nm N\'eel skyrmion emerges for $D_\uparrow$
and \emph{disappears} for $D_\downarrow$ --- controlled entirely by FE
polarization.
\item \textbf{Sign flip \& chirality reversal (trilayer).}
$D_{\uparrow\uparrow}\approx0.22$~mJ/m$^2$,
$D_{\downarrow\downarrow}\approx-0.24$~mJ/m$^2$ with
$D_{\uparrow\uparrow}\approx-D_{\downarrow\downarrow}$ and
$D_{\uparrow\uparrow}\approx D_\uparrow-D_\downarrow$; a $\sim$6~nm skyrmion whose
chirality reverses with polarization.
\item \textbf{Anisotropy caveat.} The DFT MAE ($8.57$~MJ/m$^3$) is too large
(would give $R\sim0.03$~nm); the authors adopt a moderate
$K\approx0.04$~MJ/m$^3$ to qualitatively demonstrate the FE effect.
\end{enumerate}

\section{Governing equations reimplemented}
Spin Hamiltonian (paper Eq.~5):
\begin{equation}
\mathcal{H}=-\!\sum_{\langle ij\rangle}\!J\,\hat{\mathbf m}_i\!\cdot\!\hat{\mathbf m}_j
+\!\sum_{\langle ij\rangle}\!\mathbf D_{ij}\!\cdot\!(\hat{\mathbf m}_i\!\times\!\hat{\mathbf m}_j)
-\sum_i K(\hat{\mathbf e}\!\cdot\!\hat{\mathbf m}_i)^2 ,
\end{equation}
with interfacial (N\'eel) DMI $\mathbf D_{ij}=D\,(\hat{\mathbf z}\times\hat{\mathbf u}_{ij})$.
Relaxation by overdamped LLG (projected gradient),
$\dot{\hat{\mathbf m}}=\hat{\mathbf m}\times(\hat{\mathbf m}\times\mathbf H_{\rm eff})$,
$\mathbf H_{\rm eff}=-\partial\mathcal H/\partial\hat{\mathbf m}$, normalizing each step.

Expected skyrmion radius (paper Eq.~6):
\begin{equation}
R=\Delta\,\frac{\pi D}{16\sqrt{AK}-\pi D},\qquad \Delta=\sqrt{A/K}.
\end{equation}
Verdict anchor --- the critical DMI (spiral/skyrmion instability threshold),
equivalent to the Eq.~6 denominator vanishing up to an $O(1)$ prefactor:
\begin{equation}
\boxed{\,D_c=\frac{4}{\pi}\sqrt{AK}\,}.
\end{equation}
A N\'eel skyrmion is stabilized when $|D|\gtrsim D_c$.

\section{Methods}
Discrete square lattice, $60\times60$ cells at $1$~nm spacing (matching the
paper's $60\times60$~nm supercell). Energy terms: exchange
$A/dx^2\sum(\hat{\mathbf m}_i-\hat{\mathbf m}_j)^2$, uniaxial anisotropy
$-K\sum m_z^2$, and interfacial N\'eel DMI as bond terms
$\mathbf D_{ij}\cdot(\hat{\mathbf m}_i\times\hat{\mathbf m}_j)$ with $d=D\,dx$.
A circular reversed-core seed (radius $r_0=8$~nm, smooth $180^\circ$ wall) is
relaxed for 5000 overdamped-LLG steps. \textbf{Seed handedness tracks
sign$(D)$}: $\phi=\operatorname{atan2}(Y,X)+(0\text{ if }D<0\text{ else }\pi)$,
required so the initial chirality does not oppose the DMI and spuriously
collapse the seed.

\paragraph{Fixed parameters.} $A=1.0$~pJ/m (typical Fe$_3$GeTe$_2$; the paper's
value is only in Fig.~4, not extractable from text), $K=0.04$~MJ/m$^3$ (paper's
adopted modelling value), $dx=1$~nm.

\paragraph{Diagnostics.} Berg--L\"uscher lattice topological charge
$Q=\tfrac{1}{4\pi}\sum \hat{\mathbf m}\cdot(\hat{\mathbf m}_{x+}\times\hat{\mathbf m}_{y+})$;
skyrmion diameter from reversed-core area, $2\sqrt{\text{area}/\pi}$ where $m_z<0$.
Interpreter: \texttt{/home/stevens/comfyui-env/bin/python} (numpy 2.3.5).

\section{Results}
Live re-run reproduces the saved evidence JSON to the quoted digits.
Analytic $D_c=0.2546$~mJ/m$^2$.

\begin{table}[h]
\centering
\begin{tabular}{lccccc}
\toprule
Case & $D$ (mJ/m$^2$) & vs.\ $D_c$ & Skyrmion? (this work) & $Q$ & Diameter (nm) \\
\midrule
$D_\uparrow$        & $+0.28$ & $>D_c$ & \textbf{yes} & $-0.87$ & $39.5$ \\
$D_\downarrow$      & $+0.06$ & $<D_c$ & \textbf{no}  & $\ \ 0.0$ & $0.0$ \\
$D_{\uparrow\uparrow}$   & $+0.22$ & $<D_c$$^\dagger$ & yes & $-0.90$ & $19.9$ \\
$D_{\downarrow\downarrow}$ & $-0.24$ & $|D|<D_c$$^\dagger$ & yes & $-0.99$ & $28.2$ \\
\bottomrule
\end{tabular}
\caption{This-work relaxation results. $^\dagger$The trilayer cases sit just
below the analytic $D_c$ but still relax to a metastable skyrmion from the
seed (Eq.~6 predicts a finite though small $R$ as $|D|\to D_c$); the seeded
overdamped relaxation captures the metastable branch.}
\end{table}

\begin{table}[h]
\centering
\begin{tabular}{lll}
\toprule
Quantity & This work & Paper \\
\midrule
Create/annihilate switch ($D_\uparrow$ vs $D_\downarrow$) & reproduced & core claim \\
Skyrmion present at $D_\uparrow=0.28$ & yes, $Q=-0.87$ & yes ($\sim$12 nm) \\
Skyrmion absent at $D_\downarrow=0.06$ & yes (none) & none \\
Chirality set by sign$(D)$ & reproduced (seed) & yes (trilayer) \\
Size trend (larger $|D|\to$ larger) & reproduced & $12$ vs $6$ nm \\
Absolute diameter (bilayer up) & $39.5$ nm & $\sim$12 nm \\
Critical DMI $D_c$ & $0.255$ mJ/m$^2$ & implied by Eq.~6 \\
\bottomrule
\end{tabular}
\caption{This-work vs.\ paper. Qualitative/mechanistic claims match; absolute
diameter is $\sim$2--3$\times$ too large (see gaps).}
\end{table}

\section{Comparison \& verdict}
The single most important result is that the closed-form
$D_c=(4/\pi)\sqrt{AK}=0.255$~mJ/m$^2$ lies \emph{between} the paper's two
bilayer DMI values, $D_\downarrow=0.06 < D_c < D_\uparrow=0.28$~mJ/m$^2$. This
reproduces the paper's central create/annihilate switch \textbf{analytically and
with zero fitting}: the FE polarization state, acting through the DMI it
induces, is what turns the skyrmion on and off. The seeded relaxation confirms
this dynamically ($Q=-0.87$ for $D_\uparrow$; seed collapses for $D_\downarrow$),
and the N\'eel handedness follows sign$(D)$ as the paper's trilayer
chirality-reversal claim requires. The size trend (larger $|D|$ $\to$ larger
skyrmion) is also reproduced. We therefore judge the \textbf{core mechanism
REPLICATED}, with an overall \textbf{PARTIAL} verdict driven solely by the
absolute-size mismatch below.

\section{Honest gaps \& scope}
\begin{enumerate}
\item \textbf{Exchange stiffness $A$ unknown (absolute size).} $A$ is reported
only in the paper's Fig.~4 (an image), not in the OCR/pdftotext-extractable
text. We used the typical Fe$_3$GeTe$_2$ value $A=1$~pJ/m. Since skyrmion size
$\propto\sqrt{A}$ (and $D_c\propto\sqrt{A}$), our $\sim$40~nm diameter is
$\sim$2--3$\times$ the paper's $\sim$12~nm; the \emph{threshold and trend are
correct}, the absolute scale is uncertain by $\sqrt{A_{\rm assumed}/A_{\rm true}}$.
\item \textbf{Square lattice, not the true Fe honeycomb.} The paper keeps the Fe
honeycomb; we approximate with a square lattice, which slightly changes the
geometric DMI prefactor but not the mechanism.
\item \textbf{FE polarization represented via DMI proxy.} The paper's atomistic
model has no explicit polarization field --- FE state enters only through the
DMI magnitude/sign it induces (via DFT). We correctly sweep $D$ as that proxy
rather than inventing a polarization term.
\item \textbf{No demagnetizing/dipolar field} (thin-film 2D limit, as the
paper's atomistic model also omits it).
\item \textbf{Diameter measured via reversed-core area} ($m_z<0$), a proxy for
the paper's domain-wall definition.
\end{enumerate}

\section{Reproducibility}
\begin{verbatim}
# Interpreter with numpy 2.3.5 / scipy:
/home/stevens/comfyui-env/bin/python work/skyrmion.py
# -> prints Dc and per-case {diameter, Q, skyrmion}; writes work/_sim_out.json
# Evidence copies: report/evidence/{skyrmion.py, huang2022_result.json, _sim_out.json}
\end{verbatim}
Note: no LaTeX engine on the packaging host; this \texttt{.tex} is delivered as
source and compiles off-host (\texttt{pdflatex REPORT.tex}). Extraction artifacts
(\texttt{extraction/marker.md}, \texttt{extraction/nougat.mmd}) are pdftotext
interims (marker/nougat binaries absent) --- an environment gap, not a physics gap.

\end{document}
